\chapter{Example 1.2}\label{sec:interval-example1-2}

\subsection{Problem description}\label{subsec:interval-example1-2-problem}

A café wants to plan staffing and inventory for online orders. To estimate typical demand, the manager randomly samples 50 days of order counts and summarises them in grouped form below. The goal is to estimate the mean daily orders and build a 95\% confidence interval for decision making.

Known population values: true mean (\(\mu = 11.4\)), true standard deviation (\(\sigma = 4.2\)) (used here to check the sample-based estimates).

Sample taken now: the 50 days of data are recorded in grouped form:

\begin{center}
	\begin{tabularx}{0.55\linewidth}{@{}L{0.28\linewidth}C@{}}
		\toprule
		Order count (\(x\)) & Frequency (\(f\)) \\
		\midrule
		5 & 10 \\
		9 & 15 \\
		12 & 18 \\
		20 & 7 \\
		\bottomrule
	\end{tabularx}
\end{center}

The owner wants to estimate the population mean number of daily orders and produce a 95\% confidence interval to guide staffing and stock decisions.

\subsection{Solution}\label{subsec:interval-example1-2-solution}

\subsubsection{(C1) Obtain sample mean and standard deviation}\label{subsubsec:interval-example1-2-c1}

First summarise the sample with a grouped mean to represent the population mean:

\[
\sum f x = 10(5) + 15(9) + 18(12) + 7(20) = 50 + 135 + 216 + 140 = 541
\]
\[
\bar{x} = \frac{\sum f x}{n} = \frac{541}{50} = 10.82
\]

Then measure variability using the grouped standard deviation to mirror the population spread:

\[
\sum f(x - \bar{x})^2 = 10(5 - 10.82)^2 + 15(9 - 10.82)^2 + 18(12 - 10.82)^2 + 7(20 - 10.82)^2
\]
\[
= 10(33.87) + 15(3.31) + 18(1.39) + 7(84.27) = 338.70 + 49.65 + 25.02 + 589.89 = 1003.38
\]
\[
 s = \sqrt{\frac{\sum f(x - \bar{x})^2}{n - 1}} = \sqrt{\frac{1003.38}{49}} \approx \sqrt{20.48} \approx 4.53
\]

\subsubsection{(C2) Compare sample values to the known true values}\label{subsubsec:interval-example1-2-c2}

The grouped \(\bar{x} = 10.82\) underestimates the true mean (\(\mu = 11.4\)) by about 0.6 orders, and \(s \approx 4.53\) sits slightly above the true spread (\(\sigma = 4.2\)). This illustrates how a single sample can miss the exact population values.

\subsubsection{(C3) Explain the need for estimation, based on the comparison}\label{subsubsec:interval-example1-2-c3}

Because the café can only sample a subset of days, the sample has to stand in for the ongoing population. The gaps between the sample statistics and the known population values show why we rely on estimation: samples give approximate, not exact, answers.

\subsubsection{(C5) Calculate the confidence interval (variance known)}\label{subsubsec:interval-example1-2-c5}

Compute the confidence interval for the mean at the same 95\% confidence level using the known population variance.

\[
E = z_{\alpha/2} \cdot \frac{\sigma}{\sqrt{n}} = 1.96 \cdot \frac{4.2}{\sqrt{50}} \approx 1.16
\]
\[
\mu \in 10.82 \pm 1.16 = [9.66, 11.98]
\]

\subsubsection{(C8) Calculate the confidence interval (variance unknown)}\label{subsubsec:interval-example1-2-c8}

Compute the confidence interval for the mean at the same 95\% confidence level using the sample variance.

\[
E = t_{\alpha/2,49} \cdot \frac{s}{\sqrt{n}} = t_{0.025,49} \cdot \frac{4.53}{\sqrt{50}} \approx 2.01 \cdot 0.64 \approx 1.29
\]
\[
\mu \in 10.82 \pm 1.29 = [9.53, 12.11]
\]

The t-based interval is slightly wider because it relies on the sample standard deviation.

\subsubsection{(C4) Interpret the confidence interval in real-life economic or financial terms}\label{subsubsec:interval-example1-2-c4}

For staffing and inventory budgets, the café can plan using the lower limit (about 9.5--9.7 orders) to avoid overstocking while recognising that true demand could reach the upper limit (about 12.0--12.1). The interval turns statistical uncertainty into a practical spending range.

\subsubsection{(C10) Conclude about a statistical parameter in finance and economy}\label{subsubsec:interval-example1-2-c10}

The café can conclude that the true mean daily online order count is around 10--12 orders, with a reasonable planning value near 11 orders per day. This parameter estimate informs staffing levels and inventory targets so budgets match typical demand rather than extreme days.

\vspace{6pt}
\noindent\hyperref[table-of-contents]{Back to Top}
